% 第1章 1.4.1：x^n 在 [0,1] 上的三种大小。
% (a) n=1,2,5,20 的曲线；(b) ||x^n||_inf=1，||x^n||_2=(2n+1)^{-1/2}，||x^n||_1=1/(n+1)（解析式）。
\begin{tikzpicture}[font=\figurefont\footnotesize]
\pgfplotsset{every axis/.style={width=50mm,height=38mm,axis lines=left,
  x axis line style={draw=FigureStroke,line width=.6pt},y axis line style={draw=FigureStroke,line width=.6pt},
  tick style={draw=FigureStroke},tick label style={font=\figurefont\scriptsize,text=FigureInk},
  label style={font=\figurefont\scriptsize},
  xlabel style={at={(axis description cs:1,0)},anchor=north west,yshift=1mm},
  ymajorgrids,grid style={draw=FigureGrid,line width=.4pt},
  title style={font=\figurefont\footnotesize},
  legend style={draw=none,fill=none,font=\figurefont\scriptsize,cells={anchor=west}},samples=150}}
\begin{axis}[name=A,xmin=0,xmax=1,ymin=0,ymax=1.05,xtick={0,0.5,1},ytick={0,0.5,1},xlabel={$x$},title={(a) $f_n(x)=x^n$}]
\addplot[draw=FigureBlue,line width=1.1pt,domain=0:1] {x};
\addplot[draw=FigureGreen,line width=1.1pt,domain=0:1] {x^2};
\addplot[draw=FigureYellow,line width=1.1pt,domain=0:1] {x^5};
\addplot[draw=FigurePurple,line width=1.1pt,domain=0:1,samples=300] {x^20};
\node[font=\figurefont\scriptsize,text=FigureBlue] at (axis cs:0.42,0.55) {$n=1$};
\node[font=\figurefont\scriptsize,text=FigureGreen] at (axis cs:0.62,0.30) {$2$};
\node[font=\figurefont\scriptsize,text=FigureYellow] at (axis cs:0.78,0.16) {$5$};
\node[font=\figurefont\scriptsize,text=FigurePurple] at (axis cs:0.90,0.05) {$20$};
\end{axis}
\begin{axis}[at={(A.east)},anchor=west,xshift=12mm,xmin=1,xmax=30,ymin=0,ymax=1.08,xtick={1,10,20,30},ytick={0,0.5,1},xlabel={$n$},title={(b) 三种范数},
  legend style={at={(0.98,0.80)},anchor=north east}]
\addplot[draw=FigureBlue,line width=1.1pt,domain=1:30] {1};
\addplot[draw=FigureGreen,line width=1.1pt,dashed,domain=1:30] {1/sqrt(2*x+1)};
\addplot[draw=FigureYellow,line width=1.1pt,dash dot,domain=1:30] {1/(x+1)};
\legend{$\|x^n\|_\infty$,$\|x^n\|_{2}$,$\|x^n\|_{1}$}
\end{axis}
\end{tikzpicture}
